% GPT-6 Astra Pro; effort not exposed; /root/replacement; 22 November 2026.
\section[HD-F1-R18: invariant nonsimple adjacent endpoints]{HD-F1-R18: invariant endpoint maps exclude small endomorphism sums and every rank-one semistable adjacent pair}
\label{sec:followup-rm-nonsimple-adjacent}

On the genuine real-quadratic fourfold of R07--R11,
consider actual adjacent two-cohomology objects with
\[
 \operatorname{ch}(E)=N e^{-B}v_q,\qquad
 v_q=(1-qA_1)(1-qA_2),\quad N\ne0,\quad q>0,
 \quad H^0(E)=A,\quad H^1(E)=T,
 \tag{F1R18.1}
\]
and no other cohomology, where $A,T$ are positive-rank
coherent sheaves. The following classes admit no weakly
semiregular lift for any specified coherent faithful
finite point action:
\begin{enumerate}
\item The two invariant cross-Hom groups vanish and
      $\dim\operatorname{End}_G(A)+
       \dim\operatorname{End}_G(T)\le3$.
\item The two invariant cross-Hom groups vanish,
      $A,T$ are torsion free, $|N|=1$, and the same
      endomorphism-dimension sum is at most eleven.
\item The two invariant cross-Hom groups vanish,
      $A,T$ are torsion free and slope semistable for
      one ample polarization, and $|N|=1$.
\end{enumerate}
The last exclusion has no bound on endpoint or total
endomorphism algebras. The ordinary cross-Hom groups
may be nonzero in all three statements: the required
vanishings are exactly
\[
              \operatorname{Hom}_G(A,T)=
              \operatorname{Hom}_G(T,A)=0.
 \tag{F1R18.2}
\]
Every actual connecting class, including zero, is retained.
We also prove the explicit rank and endomorphism restrictions
below for the remaining positive-rank pairs. No existence
at a permitted numerical value is asserted.

\paragraph{1. The precise equivariant endpoint injections.}
The t-exact point action induces a lift of the canonical
triangle
\[
       A\xrightarrow{i}E\xrightarrow{p}C:=T[-1]
                          \xrightarrow{\delta}A[1].
 \tag{F1R18.3}
\]
Suppose that its full degree-two equivariant trace is
injective. Put $K_E=\ker c^2_{\operatorname{ch}(E)}$.
The identity-sector forgetful comparison gives
$\operatorname{ev}_E(K_E)=0$, with all evaluations
regarded as invariant ordinary Ext classes.

The long exact Hom sequences in the actual equivariant
category contain the following maps, where the first
terms are exactly the first group in (F1R18.2):
\[
 \begin{aligned}
 \operatorname{Hom}_G(A,C[1])&\longrightarrow
 \operatorname{Ext}^2_G(A,A)
       \xrightarrow{\,i[2]\circ(-)\,}\operatorname{Hom}_G(A,E[2]),\\
 \operatorname{Hom}_G(A[1],C[2])&\longrightarrow
 \operatorname{Ext}^2_G(C,C)
       \xrightarrow{\,(-)p\,}\operatorname{Hom}_G(E,C[2]).
 \end{aligned}
 \tag{F1R18.4}
\]
In sheaf notation their first arrows are the actual
equivariant Yoneda maps
$f\mapsto\delta f$ and $f\mapsto f[2]\delta$.
Thus both displayed second arrows are injective by
(F1R18.2), even when the corresponding ordinary Hom
groups are nonzero. Naturality along $i,p$ makes
$\operatorname{ev}_A(K_E)=\operatorname{ev}_T(K_E)=0$
in equivariant Ext. Its inclusion into ordinary Ext is
injective, since it is the invariant subspace. The full
ordinary trace and the complete common annihilator in
R11 now give
\[
       \operatorname{ch}(A)=r_Ae^{-B}v_q,\qquad
       \operatorname{ch}(T)=r_Te^{-B}v_q,
                       \qquad N=r_A-r_T.
 \tag{F1R18.5}
\]
This establishes endpoint character rigidity using
exactly the invariant vanishings (F1R18.2).

\paragraph{2. The actual nonscalar endpoint boundary.}
Write
\[
 \begin{gathered}
 R_A=\operatorname{End}_G(A),\quad
 R_T=\operatorname{End}_G(T),\quad
 d_A=\dim R_A,\quad d_T=\dim R_T,\quad d=d_A+d_T,\\
 V=\operatorname{Ext}^2_G(T,A),\quad
 a=\operatorname{ext}^1_G(T,A),\quad
 b=\operatorname{ext}^1_G(A,T),\quad c=\chi_G(T,A).
 \end{gathered}
 \tag{F1R18.6}
\]
The connecting morphism has its actual endpoint-change
linear map
\[
 d_\delta:R_A\oplus R_T\longrightarrow V,\qquad
               (u,v)\longmapsto u[2]\delta-\delta v.
 \tag{F1R18.7}
\]
Here $\delta$ is written as $T\to A[2]$; equivalently
it is $C\to A[1]$. Put
$s=\dim\ker d_\delta$ and $L_\delta=\operatorname{im}d_\delta$.
Then $\dim L_\delta=d-s$.

The triangle gives the exact sequence
\[
 0\longrightarrow\operatorname{Ext}^1_G(T,A)
   \xrightarrow{\,f\mapsto ifp\,}\operatorname{End}_G(E)
   \longrightarrow\ker d_\delta\longrightarrow0,
 \tag{F1R18.8}
\]
where in its first arrow $f$ is viewed as $C\to A$,
and the last arrow takes the actual induced maps on
$H^0(E),H^1(E)$. To verify its exactness without a
filtration-preservation assumption, apply the two Hom
sequences. First
$\operatorname{Hom}_G(E,C)=R_Tp$ because
$\operatorname{Hom}(A,C)=\operatorname{Hom}(A[1],C)=0$.
Also $\operatorname{Hom}(E,C[-1])=0$. The sequence
for $\operatorname{Hom}(-,A)$ begins
\[
 0\longrightarrow\operatorname{Ext}^1_G(T,A)
   \longrightarrow\operatorname{Hom}_G(E,A)
   \longrightarrow R_A\xrightarrow{\,u\mapsto u[2]\delta\,}V.
\]
The image of $\operatorname{End}_G(E)\to R_T$ is
therefore precisely the set of $v$ for which
$\delta v$ belongs to $R_A[2]\delta$; the fibre over
each compatible pair has kernel the first group in
(F1R18.8). This proves the sequence and
\[
                  e_0:=\operatorname{ext}^0_G(E,E)=a+s.
 \tag{F1R18.9}
\]

The tangent map $J(v)=i[1]vp$ has kernel exactly
$L_\delta$, so it gives an injection
\[
       V/L_\delta\lhook\joinrel\longrightarrow
                              \operatorname{Ext}^1_G(E,E).
 \tag{F1R18.10}
\]
Indeed the kernel of
$V\xrightarrow{(-)p}\operatorname{Hom}_G(E,A[1])$
is $R_A[2]\delta$. The kernel of the next map induced
by $i[1]$ is the image of
$\operatorname{Hom}_G(E,C)=R_Tp$ under composition
with $\delta$. Taking its inverse image gives exactly
$R_A[2]\delta+\delta R_T=L_\delta$. These statements
hold also when $\delta=0$.

The eight-dimensional characteristic tangent image is
disjoint from (F1R18.10). For an equality
$\operatorname{ev}_E^1(x)=J(v)$, compose with $i$.
The zero composite $pi$ and naturality give
$i[1]\operatorname{ev}_A^1(x)=0$.
The actual map
\[
 \operatorname{Ext}^1_G(A,A)\longrightarrow
                    \operatorname{Hom}_G(A,E[1])
 \tag{F1R18.11}
\]
is injective because its preceding term is
$\operatorname{Hom}_G(A,C)=\operatorname{Ext}^{-1}_G(A,T)=0$.
Equation (F1R18.5) makes the full degree-one character
map of $A$ injective on $HH^1(X)$. Hence $x=0$.
It follows that
\[
 e_1:=\operatorname{ext}^1_G(E,E)
                   \ge8+\dim V-d+s.
 \tag{F1R18.12}
\]

\paragraph{3. The exact integral endpoint restriction.}
The second vanishing in (F1R18.2) makes all negative
cross-Hom groups between $A$ and $C$ vanish in the
equivariant category. The two Hom sequences then give
$\operatorname{Ext}^j_G(E,E)=0$ for $j<0$; ordinary
negative Ext vanishing is not needed. Invariant CY4
duality, the full rank-eighteen trace and the cross
degree-zero and degree-four vanishings give
\[
 \dim V=c+a+b,\quad e_2=18,\quad
 k:=\chi_G(E,E)=2e_0-2e_1+18.
 \tag{F1R18.13}
\]
Using (F1R18.9) and (F1R18.12), the parameter $s$
cancels exactly and yields
\[
                            c+b+\frac{k}{2}\le d+1.
 \tag{F1R18.14}
\]
The already proved induction/forgetful Euler identity,
polarized on the three actual objects $A,T,A\oplus T$,
and (F1R18.5) give
\[
 c=\frac{4Dq^2r_A r_T}{|G|}>0,
 \qquad k=\frac{4Dq^2(r_A-r_T)^2}{|G|}>0,
 \qquad c(r_A-r_T)^2=k r_A r_T.
 \tag{F1R18.15}
\]
Thus $c$ is a positive integer and $k$ is a positive
even integer.

Put $g=\gcd(r_A,r_T)$, $p=r_A/g$, $q_0=r_T/g$,
so $p,q_0$ are positive coprime unequal integers.
Since $\gcd(pq_0,(p-q_0)^2)=1$, (F1R18.15) gives
a positive integer $t$ such that
\[
 \begin{gathered}
 k=t(p-q_0)^2,\qquad c=tpq_0,\qquad t(p-q_0)^2\in2\mathbf Z,\\
                    d+1\ge\frac{t(p^2+q_0^2)}2+b.
 \end{gathered}
 \tag{F1R18.16}
\]
The last main term is at least five. If $p-q_0$ is
odd, $t\ge2$ and $p^2+q_0^2\ge5$. If it is even,
coprimality makes both numbers odd, so their unequal
positive values give $p^2+q_0^2\ge10$ and $t\ge1$.
Consequently
\[
                                  d\ge4+b.
 \tag{F1R18.17}
\]
This proves the first exclusion. At $d=4$ it forces
$b=0$, and the only rank-ratio possibilities, up to
interchanging the endpoints, are $(p,q_0,t)=(2,1,2)$
and $(3,1,1)$. These are necessary profiles only.

For torsion-free endpoints, S-22 and (F1R18.5) give
$r_A,r_T\ge2$. When $|N|=1$, their gcd is one and
$|p-q_0|=1$, so $t\ge2$ and
$p^2+q_0^2\ge3^2+2^2=13$. Hence
\[
                  |N|=1,\ A,T\text{ torsion free}
                  \quad\Longrightarrow\quad d\ge12+b.
 \tag{F1R18.18}
\]
This proves the second exclusion without a semistability
assumption.

\paragraph{4. An actual semistable-filtration bound for each endpoint.}
We next prove a stronger ordinary bound for every actual
slope-semistable torsion-free sheaf $Q$ with
\[
                    \operatorname{ch}(Q)=r e^{-B}v_q.
 \tag{F1R18.19}
\]
Its rank satisfies $r\ge2$ by S-22. Choose a saturated
slope Jordan--H\"older filtration with stable torsion-free
quotients $F_1,\ldots,F_m$, all of slope $\mu(Q)$.
This is an actual coherent-sheaf filtration: when $Q$
is not stable, take a same-slope subsheaf of least
positive rank and saturate it. Saturation can only
increase slope, while semistability of $Q$ prevents
an increase. Least rank makes the resulting subsheaf
stable. The torsion-free quotient is semistable of the
same slope, because the preimage of a higher-slope
subsheaf would destabilize $Q$. Induction on rank
constructs the asserted finite filtration.

Here are the exact Hom maps controlling its endomorphisms.
For $0\to Q_{j-1}\to Q_j\to F_j\to0$, the sequences
\[
 \begin{split}
 0&\longrightarrow\operatorname{Hom}(Q,Q_{j-1})
  \longrightarrow\operatorname{Hom}(Q,Q_j)
  \longrightarrow\operatorname{Hom}(Q,F_j),\\
 0&\longrightarrow\operatorname{Hom}(F_i,F_j)
  \longrightarrow\operatorname{Hom}(Q_i,F_j)
  \longrightarrow\operatorname{Hom}(Q_{i-1},F_j)
 \end{split}
 \tag{F1R18.20}
\]
give a finite filtration of the Hom space with successive
quotients embedding in the displayed stable-factor Hom
spaces. In particular
\[
              \dim\operatorname{End}(Q)
                   \le\sum_{i,j}\dim\operatorname{Hom}(F_i,F_j).
 \tag{F1R18.21}
\]
No endomorphism is presumed to preserve the chosen
Jordan--H\"older filtration; (F1R18.20) provides the
required actual maps in both variables.

A nonzero morphism between stable torsion-free sheaves
of the same slope has equal endpoint ranks and full
generic rank, by the image-and-kernel argument of R16.
Each reflexive hull remains slope stable: intersect a
proper positive-rank subsheaf of $F_i^{**}$ with $F_i$;
the quotient has codimension at least two, so its rank
and first Chern class are unchanged and stability of
$F_i$ applies. The double-dual morphism has a nonzero
determinant section of a line bundle of degree zero.
Its effective zero divisor must be empty, since the
chosen polarization is ample. It is therefore an
isomorphism outside a set of codimension at least two,
and extends to an isomorphism of reflexive sheaves by
their actual identities $F_i^{**}=j_*(F_i^{**}|_U)$ on
that common big open set. Stable reflexive sheaves have
only scalar endomorphisms, by the proof in R16. Finally
double dual induces an injective map
\[
 \operatorname{Hom}(F_i,F_j)\hookrightarrow
                   \operatorname{Hom}(F_i^{**},F_j^{**}),
 \tag{F1R18.22}
\]
because the target $F_j$ is torsion free. It follows
that each stable-factor Hom dimension is at most one,
and is zero unless the two actual reflexive hulls are
isomorphic.

Group the factors by their reflexive-hull isomorphism
classes, with multiplicities $m_\alpha$ and hull ranks
$\rho_\alpha$. Equations (F1R18.21)--(F1R18.22) give
\[
       \dim\operatorname{End}(Q)\le\sum_\alpha m_\alpha^2,
        \qquad r=\sum_\alpha m_\alpha\rho_\alpha.
 \tag{F1R18.23}
\]
If one $\rho_\alpha\ge2$, the number of factors is
at most $r-1$, so the right side is at most $(r-1)^2$.
If every hull is a line bundle and at least two hull
types occur, the positive multiplicities sum to $r$
and their squared sum is at most $(r-1)^2+1$.

The sole remaining case is that all $r$ factors have
one actual line-bundle hull $L$. Let $\ell=c_1(L)$.
The quotients $L/F_i$ have codimension at least two,
and their codimension-two character classes $Z_i$ are
effective surface cycles. Actual additivity along the
filtration gives
\[
      c_1(Q)=r\ell=-rB,\qquad
      \operatorname{ch}_2(Q)
          =\frac r2\ell^2-\sum_i Z_i
          =\frac r2\ell^2-rq(A_1+A_2).
 \tag{F1R18.24}
\]
Thus $\sum_i Z_i=rq(A_1+A_2)$ is a nonzero effective
cycle in the plane excluded by the exact Gauss-map and
rational-eigenvector proof of S-22. This case is
impossible. We have proved the unconditional bound
\[
       \dim\operatorname{End}_X(Q)\le(r-1)^2+1
                                      =r^2-2r+2
 \tag{F1R18.25}
\]
for every actual sheaf in (F1R18.19).

\paragraph{5. The nonsimple semistable adjacent exclusion.}
For the two slope-semistable torsion-free endpoints,
(F1R18.5) forces equality of their normalized first
Chern classes, hence of slopes for the chosen ample
polarization. Put
$R=r_A^2+r_T^2$ and $S=r_A+r_T$. Taking invariant
endomorphisms is passage to a subspace, so
(F1R18.25) gives
\[
                              d+1\le R-2S+5<R,
 \tag{F1R18.26}
\]
where $r_A,r_T\ge2$ makes the strict inequality
immediate. On the other hand, (F1R18.14)--(F1R18.15)
give the exact bound
\[
                         \frac{kR}{2N^2}+b\le d+1.
 \tag{F1R18.27}
\]
Thus $k<2N^2$. Since $k$ is a positive even integer,
\[
                             2\le k\le2N^2-2.
 \tag{F1R18.28}
\]
There is no such integer when $|N|=1$, proving the
third exclusion with no endomorphism bound. The same
calculation shows that the minimal profile
$(e_0,e_1,e_2)=(3,8,18)$, for which $k=8$, cannot
arise from these adjacent semistable endpoints when
$|N|\le2$.

\paragraph{Framework and exact remaining scope.}
The invariant endpoint injections, the nonscalar
boundary $d_\delta$, its actual stabilizer kernel and
tangent quotient, and the independent point-action
map are all written above. The character operations
are actual $K_0$ additivity along the canonical
triangle and the saturated coherent-sheaf filtration.
The effective-cycle contradiction uses the actual
rank-one quotients $L/F_i$, not a proposed decomposition
of a numerical character. The reflexive-Hom map is
double dual followed by the specified big-open
extension. The equivariant Euler pairing retains the
same published induction/forgetful and Hochschild
Riemann--Roch maps as R16--R17.
The precise big-open extension identity used in
(F1R18.22) is the restriction/pushforward inverse of
Stacks Project Lemma~31.13.12, in Section~31.13
(Tag~0AVT); its complete proof was reopened. The
saturated slope filtration and the stable-factor Hom
bound are proved directly in paragraph~4, so neither
a preservation assertion for endomorphisms nor a
numerical rank-square estimate is imported.

The exclusions do not address arbitrary adjacent
objects with nonzero invariant cross-Hom groups. The
all-endomorphism semistable exclusion is at virtual
rank one in absolute value; (F1R18.28) is a proved
necessary bound, not a construction or general
nonexistence result in larger virtual rank. No
unrefereed theorem is a proof premise.
